\documentclass[11pt]{article}
\ifdefined\FONTS\else\def\FONTS{ot1}\fi
\def\fontsT{t1}\def\fontsL{lm}\def\fontsP{ptm}
\ifx\FONTS\fontsT \usepackage[T1]{fontenc}\fi
\ifx\FONTS\fontsL \usepackage[T1]{fontenc}\usepackage{lmodern}\fi
\ifx\FONTS\fontsP \usepackage[T1]{fontenc}\usepackage{mathptmx}\usepackage[scaled=0.9]{helvet}\usepackage{courier}\fi
\usepackage[utf8]{inputenc}

\usepackage{amsmath,amssymb,amsthm}
\newtheorem{theorem}{Theorem}
\title{On the Convergence of Simple Iterative Methods}
\author{Ada Example \and Brook Sample}
\date{March 2004}
\begin{document}
\maketitle
\begin{abstract}
We study the convergence of fixed point iterations for contractive maps on complete metric spaces.
The main result gives an explicit bound on the number of steps needed to reach a given accuracy,
and we show that the bound is sharp for a natural family of affine maps.
\end{abstract}
\section{Introduction}
Let $(X,d)$ be a complete metric space and let $T\colon X\to X$ satisfy
$d(Tx,Ty)\le q\,d(x,y)$ for all $x,y\in X$, where $0\le q<1$. The classical theorem of Banach states
that the sequence $x_{n+1}=Tx_n$ converges to the unique fixed point $x^\ast$ of $T$ for every starting value $x_0$.
In this note we revisit the classical argument and make every constant explicit.
\begin{theorem}
For every $\varepsilon>0$ the iteration satisfies $d(x_n,x^\ast)<\varepsilon$ whenever
\begin{equation}
n \ge \frac{\log\varepsilon - \log d(x_0,x_1) + \log(1-q)}{\log q}.
\end{equation}
\end{theorem}
\begin{proof}
By the triangle inequality and the contraction property,
\[
d(x_n,x^\ast)\le\sum_{k=n}^{\infty} d(x_k,x_{k+1})\le \sum_{k=n}^{\infty} q^k d(x_0,x_1)=\frac{q^n}{1-q}\,d(x_0,x_1).
\]
Taking logarithms gives the claim.
\end{proof}
\section{Examples}
For the affine map $T(x)=\alpha x+\beta$ on $\mathbb{R}$ with $|\alpha|<1$ the fixed point is
$x^\ast=\beta/(1-\alpha)$, and the error after $n$ steps equals $|\alpha|^n|x_0-x^\ast|$.
The integral $\int_0^1 e^{-x^2}\,dx\approx 0.7468$ and the sum $\sum_{k\ge1} k^{-2}=\pi^2/6$ appear in the
second example, together with the Greek letters $\alpha,\beta,\gamma,\delta,\lambda,\mu,\sigma,\omega$ and $\Gamma,\Delta,\Omega$.
We also use the sets $\mathbb{N}\subset\mathbb{Z}\subset\mathbb{Q}\subset\mathbb{R}$ and the relation $a\approx b$.
\section{Conclusion}
The estimate is simple, but it is often all that is needed in practice. Further work will consider
maps that are only contractive on average.
\end{document}
